\section{Motivation}\label{sec:id-motivation}
\owner{Group}
\pages{1}

Students are more likely to persist with a task they find relevant, that they believe they can complete, and that gives them some control. The app draws on two frameworks to support this: the ARCS model \cite{Keller2010} and self-determination theory \cite{Deci2000}.

\subsection{ARCS Model}
The ARCS model names four conditions for motivated learning: attention, relevance, confidence and satisfaction \cite{Keller2010}. The app addresses each one as follows.

\begin{description}
    \item[Attention.] The predicted-step questions ask the student to commit to an answer before the explanation, which makes each step a small question rather than a passive read.
    \item[Relevance.] The introduction to each lesson states how the lesson connects to the lecture and to the search performance that the student will meet in later modules.
    \item[Confidence.] Success criteria are stated at the start of each lesson, and the graduated hints let a student succeed with help before attempting without it.
    \item[Satisfaction.] Each completed objective is marked in the progress record, so the student sees the result of their work.
\end{description}

\subsection{Self-Determination Theory}
Self-determination theory identifies three basic psychological needs: autonomy, competence and relatedness \cite{Deci2000}. The app addresses two of them directly. Autonomy is supported because the student controls the pace of every trace and can repeat any lesson. Competence is supported because progress is shown as an accurate record of what the student has practised, and because the hints are visibly withdrawn as competence grows. Relatedness is not addressed by the prototype. The app is used alone, and the group has not designed any social feature. This is recorded as a limitation in the SE document.

\subsection{Gamification}
The app uses only visible progress as a game element. It has no points, no leaderboard and no timed challenge. Timed challenges can reward speed over understanding, and a student who is anxious about a timer may stop trying. Progress bars were chosen because they reflect the mastery standard in Section~\ref{sec:id-progress} and say nothing about how other students are doing.
